% ============================================================
%  STILE DEL FORMULARIO DI METODI STATISTICI PER L'INFORMAZIONE
% ============================================================
% Si carica dopo preamble/mspi.tex, da cui prende colori e notazione (\compl, \card,
% \vett, \matr, \decidi, \dd, \abs...). Aggiunge:
%   \argomento{Titolo}{capitoli}        sezione numerata (1, 2, ...); con [R] la sezione
%                                       dei richiami di matematica (voci R.1, R.2, ...)
%   \voce{etichetta}{Nome}{nota}{formula}{varianti}
%                                       una voce della prima parte: la formula va scritta
%                                       come matematica (è composta in display style)
%   \variante{quando}{formula}          una variante, dentro l'ultimo argomento di \voce
%   \oltre                              etichetta «oltre le slide» per le formule aggiunte
%   \scheda{etichetta}{righe}           la scheda d'uso della seconda parte, con le righe
%                                       \usoquando, \usocome, \usovarianti, \usoattenzione,
%                                       \usoesempio
%   ambiente scelta{Titolo}             il riquadro «come scegliere» di un argomento
% Il numero di ogni voce porta alla sua scheda d'uso, e la scheda riporta alla voce.

% ------------------------------------------------------------
% Colori
% ------------------------------------------------------------
\colorlet{fcolore}{mspiform!80!black}   % voci delle formule, come i riquadri «Da ricordare»
\colorlet{fsfondo}{mspiform!5}          % sfondo alternato delle voci
\colorlet{ucolore}{mspiteo!85!black}    % schede d'uso, come i teoremi
\colorlet{acolore}{mspiB!85!black}      % righe «Attenzione»

% ------------------------------------------------------------
% Argomenti e numerazione
% ------------------------------------------------------------
\newcounter{argomento}
\newcounter{voce}
\def\numargomento{}
\renewcommand{\thevoce}{\numargomento.\arabic{voce}}
% Nome dell'ancora del PDF: deve essere diverso per ogni voce (voce.2.4, voce.R.1...), altrimenti
% i collegamenti alle voci di argomenti diversi con lo stesso numero finirebbero tutti sulla prima
\renewcommand{\theHvoce}{\numargomento.\arabic{voce}}
\newif\ifparteuso                         % vero nella seconda parte (come si usano)

% \argomento[numero speciale]{Titolo}{nota sotto il titolo}[titolo breve per il piè di pagina]
\NewDocumentCommand{\argomento}{o m m o}{%
  \IfNoValueTF{#1}%
    {\stepcounter{argomento}\xdef\numargomento{\arabic{argomento}}\colorlet{fcolore}{mspiform!80!black}%
     \colorlet{fsfondo}{mspiform!5}}%
    {\xdef\numargomento{#1}\colorlet{fcolore}{mspiric!85!black}\colorlet{fsfondo}{mspiric!5}}%
  \setcounter{voce}{0}%
  \section{\numargomento\enspace #2}%
  \IfNoValueTF{#4}{\gdef\argomentobreve{#2}}{\gdef\argomentobreve{#4}}%
  \ifparteuso\gdef\currentchapter{Come si usa: \argomentobreve}\else\gdef\currentchapter{\argomentobreve}\fi
  \vspace*{-4mm}{\Lato\small\color{mspigray}#3}\par\nobreak\vspace{2mm}%
}

% ------------------------------------------------------------
% Le voci del formulario
% ------------------------------------------------------------
% Numero in un rettangolino colorato
\newcommand{\badge}[2][fcolore]{%
  {\setlength{\fboxsep}{1.5pt}\colorbox{#1}{\Lato\bfseries\footnotesize\color{white}\strut\,#2\,}}}

\newtcolorbox{vocebox}[1][]{%
  enhanced, breakable=false, sidebyside, sidebyside align=top seam, sidebyside gap=5mm,
  lefthand width=52mm, boxrule=0pt, frame hidden, arc=0pt, colback=white,
  left=3mm, right=3mm, top=2.2mm, bottom=2.2mm, before skip=0pt, after skip=0pt,
  borderline south={0.5pt}{0pt}{mspigray!30},
  segmentation style={solid, line width=0.5pt, mspigray!30},
  before upper={\setlength{\parskip}{1pt}\raggedright}, before lower={\setlength{\parskip}{0pt}\raggedright},
  #1}

% \voce{etichetta}{Nome}{nota sotto il nome}{formula}{varianti}
\NewDocumentCommand{\voce}{m m m +m +m}{%
  \expandafter\gdef\csname fnome@#1\endcsname{#2}%
  \ifodd\value{voce}\def\vocesfondo{fsfondo}\else\def\vocesfondo{white}\fi
  \begin{vocebox}[colback=\vocesfondo]
    \refstepcounter{voce}\label{f:#1}%
    \hyperref[u:#1]{\badge{\thevoce}}\hfill
      {\Lato\scriptsize\hyperref[u:#1]{\color{mspigray}come si usa: p.~\pageref*{u:#1}}}\par
    \vspace{2pt}{\bfseries #2}\par
    \ifx\relax#3\relax\else{\footnotesize\color{mspigray}#3\par}\fi
  \tcblower
    $\displaystyle #4$%
    #5%
  \end{vocebox}%
}

% \variante{quando}{formula}: una riga sotto la formula principale
\newcommand{\variante}[2]{%
  \par\vspace{3.5pt}\noindent\hangindent=1.2em\hangafter=1
  {\color{fcolore}\footnotesize$\blacktriangleright$}\enspace
  \ifx\relax#1\relax\else{\small\itshape #1:}\enspace\fi
  $\displaystyle #2$}

% Etichetta per le formule che non sono nelle slide (generalizzazioni comode, sempre vere)
\newcommand{\oltre}{%
  {\setlength{\fboxsep}{0.8pt}\fcolorbox{acolore}{white}{\Lato\upshape\scriptsize\color{acolore}\,oltre le slide\,}}}

% ------------------------------------------------------------
% Le schede d'uso
% ------------------------------------------------------------
\newtcolorbox{schedabox}[2]{%
  enhanced, breakable=false, colback=white, colframe=mspiteo!45, boxrule=0.5pt, arc=1.5mm,
  left=3mm, right=3mm, top=1.5mm, bottom=2.5mm, before skip=2.5mm, after skip=2.5mm,
  colbacktitle=mspiteo!9, coltitle=ucolore, fonttitle=\bfseries, toptitle=0.9mm, bottomtitle=0.9mm,
  titlerule=0.5pt, titlerule style={mspiteo!30},
  title={{\setlength{\fboxsep}{1.4pt}\colorbox{ucolore}{\Lato\bfseries\footnotesize\color{white}\strut\,\ref*{f:#1}\,}}%
         \enspace #2\hfill{\Lato\footnotesize\mdseries\hyperref[f:#1]{formula a p.~\pageref*{f:#1}\enspace$\uparrow$}}},
  before upper={\setlength{\parskip}{0pt}},
}

\NewDocumentCommand{\scheda}{m +m}{%
  \ifcsname fnome@#1\endcsname\else\PackageError{formulario}{Voce #1 inesistente}{}\fi
  \begin{schedabox}{#1}{\csname fnome@#1\endcsname}%
    \phantomsection\label{u:#1}%
    #2%
  \end{schedabox}}

% Una riga della scheda: etichetta a sinistra, testo a destra
\newlength{\usoetichetta}
\setlength{\usoetichetta}{21mm}
\newcommand{\rigauso}[3]{%
  \par\vspace{2.5pt}\noindent
  % \noindent prima del colore: il riferimento delle minipage [t] deve essere la prima riga di testo
  \begin{minipage}[t]{\usoetichetta}\noindent{\Lato\bfseries\small\color{#1}#2}\end{minipage}%
  \begin{minipage}[t]{\dimexpr\linewidth-\usoetichetta\relax}\setlength{\parskip}{3pt}\raggedright\noindent#3%
  \end{minipage}\par}
\newcommand{\usoquando}[1]{\rigauso{ucolore}{Quando}{#1}}
\newcommand{\usocome}[1]{\rigauso{ucolore}{Come}{#1}}
\newcommand{\usovarianti}[1]{\rigauso{fcolore}{Varianti}{#1}}
\newcommand{\usoattenzione}[1]{\rigauso{acolore}{Attenzione}{#1}}
\newcommand{\usoesempio}[1]{\rigauso{mspigray!70!black}{Esempio}{#1}}

% Rimando a un'altra voce, cliccabile: \vv{etichetta} scrive il suo numero
\newcommand{\vv}[1]{\hyperref[f:#1]{\textcolor{mspiform!80!black}{\Lato\bfseries\ref*{f:#1}}}}

% Riquadro «come scegliere la formula», all'inizio di ogni argomento della seconda parte
\newtcolorbox{scelta}[1]{%
  enhanced, breakable=false, colback=primary!4, colframe=primary!55, boxrule=0.5pt, arc=1.5mm,
  left=3mm, right=3mm, top=2mm, bottom=2.5mm, before skip=1mm, after skip=4mm,
  fonttitle=\bfseries, colbacktitle=primary!55, coltitle=white, title={#1},
  before upper={\setlength{\parskip}{4pt}},
}
